% Nature Reviews Neuroscience: source-inspired original neutral demonstration.
% No real article text, scientific data, author list, DOI, or original figures.
\PassOptionsToPackage{no-math}{fontspec}
\documentclass[UTF8,fontset=none,scheme=plain]{ctexart}
\usepackage[paperwidth=595.276bp,paperheight=790.866bp,left=39.685bp,right=34.016bp,top=114bp,bottom=43bp,headheight=20bp,headsep=42bp,footskip=24bp]{geometry}
% Shared dependencies and neutral material. Compile with XeLaTeX.
% No publisher logos, paper text, data, figures, or download metadata are included.
\usepackage{amsmath,amssymb,graphicx,xcolor,tikz,fancyhdr,needspace,xurl,hyperref,balance,etoolbox}
% The legacy balance package omits the first split box's final baseline depth.
% Reserve that measured depth rather than hiding a CJK overfull-box warning.
\makeatletter
\patchcmd{\@BAlancecol}{\setbox\@leftcolumn=\vbox to \@tempdima}{\advance\@tempdima by \dp1\relax\setbox\@leftcolumn=\vbox to \@tempdima}{}{\PackageWarning{cns-template}{Balance depth patch unavailable; inspect final columns}}
\makeatother
\usetikzlibrary{arrows.meta,positioning}
\usepackage[protrusion=true,expansion=false]{microtype}
\IfFontExistsTF{Noto Serif CJK SC}{%
  \setCJKmainfont{Noto Serif CJK SC}%
}{%
  \setCJKmainfont[BoldFont=FandolSong-Bold.otf,ItalicFont=FandolKai-Regular.otf]{FandolSong-Regular.otf}%
}
\IfFontExistsTF{Noto Sans CJK SC}{%
  \setCJKsansfont{Noto Sans CJK SC}%
}{%
  \setCJKsansfont[BoldFont=FandolHei-Bold.otf]{FandolHei-Regular.otf}%
}
\setmainfont{TeX Gyre Termes}\setsansfont{TeX Gyre Heros}
\setmonofont{TeX Gyre Cursor}
\definecolor{JournalBlue}{HTML}{0097C9}
\definecolor{JournalRed}{HTML}{B83443}
\definecolor{ScienceRed}{HTML}{CD1021}
\definecolor{CiteBlue}{HTML}{0000EE}
\hypersetup{colorlinks=true,linkcolor=JournalBlue,urlcolor=JournalBlue,
  pdftitle={CNS论文汉化：中性版式示例},pdfauthor={示例模板；非期刊官方模板},
  pdfsubject={Neutral source-inspired Chinese layout demonstration; no real research claims}}
\urlstyle{same}\raggedbottom
\setlength{\parindent}{1em}\setlength{\parskip}{0pt}
\setlength{\emergencystretch}{2em}\tolerance=2200
\XeTeXlinebreaklocale "zh"\XeTeXlinebreakskip=0pt plus .1pt
\pagestyle{fancy}\fancyhf{}
\renewcommand{\headrulewidth}{0pt}\renewcommand{\footrulewidth}{0pt}
\newcommand{\PaperTitle}{中文题名占位：从组织结构到证据解释}
\newcommand{\SampleAuthors}{Author A\textsuperscript{1}, Author B\textsuperscript{2,*}}
\newcommand{\SampleAffiliations}{\textsuperscript{1}作者单位一（示例占位）。\quad
  \textsuperscript{2}作者单位二（示例占位）。\par
  *通信作者信息待填。DOI、收稿和出版日期均应从目标原文核实。}
\newcommand{\SampleAbstract}{本段是中性摘要占位文字，不陈述任何真实研究发现。它演示中文题名、作者信息、摘要和双栏正文之间的层级关系。实际使用时，应完整替换为目标论文的忠实中文译文，并保留研究对象、方法、结果和结论的限定条件。图注、参考文献、作者声明和来源信息均属于独立的覆盖项。模板中的图形仅说明排版结构，不表示生物学机制、实验结果或统计关系。}
\newcommand{\DemoCaption}{\textbf{图1. 中性证据链示意（占位图，无实验数据）。}
  \textbf{A}，输入节点表示需要核对的原文材料。
  \textbf{B}，处理节点表示逐段转换及复核。
  \textbf{C}，输出节点表示正文、图注和书目等独立交付部分。
  箭头仅表示本示例的阅读顺序，不代表因果关系、效应方向或数量变化。
  该图由模板自行绘制，所有节点均为文字占位；无样本量、误差线、统计检验或比例尺。
  实际论文图注应逐项保留原稿报告的样本量、单位、检验方法、限定及图形署名。}
\newcommand{\DemoVector}{%
 \resizebox{\linewidth}{!}{%
 \begin{tikzpicture}[font=\sffamily\small,node distance=18mm,
   box/.style={draw=black!70,rounded corners=2pt,fill=black!3,minimum width=26mm,minimum height=16mm,align=center},
   flow/.style={-{Stealth[length=2.5mm]},line width=.7pt}]
  \node[box] (a) {\textbf{A}\quad 输入\\原文材料};
  \node[box,right=of a] (b) {\textbf{B}\quad 处理\\转换与复核};
  \node[box,right=of b] (c) {\textbf{C}\quad 输出\\中文译稿};
  \draw[flow] (a)--(b);\draw[flow] (b)--(c);
  \node[below=11mm of a,align=center] {标题与摘要\\正文与图表};
  \node[below=11mm of b,align=center] {保留限定\\检查单位};
  \node[below=11mm of c,align=center] {可追溯引文\\明确原文信息};
  \node[below=29mm of b,font=\sffamily\footnotesize,text=black!65] {示意占位；不包含真实实验数据};
 \end{tikzpicture}}%
}
\newcommand{\SampleBibliography}{%
 \section*{参考文献（示例占位）}
 {\fontsize{7.6bp}{10bp}\selectfont\setlength{\parindent}{0pt}\interlinepenalty=10000
 \hypertarget{ref:1}{}1. Author A. (20XX). Example reference title. \textit{Journal placeholder}. 卷、页码与DOI待填。此条仅供演示引文跳转。\par\smallskip
 \hypertarget{ref:2}{}2. Author B. (20YY). Another example reference title. \textit{Journal placeholder}. 出版信息待填；不要把占位条目带入正式译稿。\par}
}
\newcommand{\SampleSourceInfo}{%
 \section*{原文信息}
 以下字段全部为示例占位：作者；英文原标题；期刊；年份；卷期；页码或文章号；DOI或来源链接。实际使用时必须逐项替换并核对。\par
 本文件是\JournalName 风格的非官方中文译稿样板，依据一次已核验文章版式进行结构化改编，不代表该期刊所有年代或文体的统一标准，也不保证像素级复刻。无真实论文内容或期刊标识图像。\par
}
% Deliberately neutral demo prose, not translated from a published paper.
\newcommand{\SampleOpening}{%
 \section*{引言：示例正文与引文}
 这是一段用于检查双栏阅读节奏的中性示例文字。段落中的术语、缩写和标点需要在实际项目中与源文保持一致。中文译稿应呈现原论文的完整论述，而不是把多段证据压缩成一段摘要。这里的引文只链接到本文件末尾明确标为占位的参考文献\SampleCiteOne。\par
 若原句仅表示相关、可能或提示，译文也应保留相同的证据强度。数字、单位、正负号、上下标和范围端点应逐项核对。原文存在内部矛盾时，应分别保留相应位置的陈述，再用简短译注说明；不能根据常识把不同的数字静默统一。\par
 \subsection*{内容层级与跨页衔接}
 此小标题演示二级层级。实际论文可能采用不同的小标题密度，模板不会替代对原文结构的识别。双栏文章需按左栏、右栏和跨页续句的顺序读取。若段落在页末中断，应在完成整句核对后接入下一页，避免把图中文字误当正文，或重复翻译同一段。\par
 图表具有独立的信息范围。正文引用到某幅图，并不意味着其图注可以省略。图注中的每个分面、比较对象、时间点、样本量及统计说明都应对应源文。英文图内标签可以保留，并在图旁或文末提供中文释义；实际采用哪种安排，需要依据图形尺寸和译文长度判断\SampleCiteTwo。\par
 \subsection*{来源与可维护数据}
 本段继续演示正文密度。可编辑文本和图注数据应与最终版面分开维护，以便修订时准确定位变化。若只调整版式，应能证明正文、图注和书目内容没有被意外改写。更换字体或纸张后，须重新检查断行、图注与图片的相对位置以及页脚可读性。\par
 文内引用与书目之间应有明确对应关系。这里的两条演示书目没有真实DOI，也没有虚构的实验信息。使用模板时可增删条目，但必须同时更新编号、链接目标和引用位置。页码只表示中文译稿自己的分页，不用于冒充原出版页码。\par
 \subsection*{结果叙述的演示占位}
 本节的文字仅用于呈现长段落在双栏页面中的流动，不报告真实结果。实际翻译应保留观察对象、实验条件和结局指标之间的关系。例如，原文分别描述体外实验与动物实验时，不能把两类结果合并为同一层面的结论；研究使用的对照、时间顺序和分析范围也必须保留。图表中的比较范围不应扩大为对整个人群或全部疾病阶段的判断。\par
 论述中的转折与限制同样属于正文。若研究先提出一种解释，再讨论另一种可能，译文应沿用这一推进顺序。作者明确写出的不确定性应能在中文中被识别；如果尚未确定某个现象的方向或机制，就不能改成已经建立的因果关系。对于术语近似但研究层级不同的概念，可在首次出现时保留缩写，后文保持统一。\par
 \subsection*{方法信息与可复核性}
 方法部分往往包含密集数字和专有名词。翻译时须区分浓度、总量、频率和处理时长，不能因为中文句式调整而改变参数的归属。试剂名称、软件版本和数据访问标识通常保留原样；上下限、排除规则和显著性标准则需要逐字符核对。这里没有提供任何真实实验参数，避免把模板占位误认为可以直接执行的方法。\par
 表格与补充材料的引用也需要保留原层级。正文对补充图的指向并不意味着模板中已经附有该补充图；若源文件只包含主文，应如实标注所交付的范围。图表编号、附录顺序和参考文献数量都应由目标文章决定，本示例中的结构不构成固定配额，也不要求所有论文采用同样的章节。\par
}
\newcommand{\SampleClosing}{%
 \section*{讨论：版式适配与范围}
 上方矢量图是模板原创的排版占位图，不来自任何论文。图形与完整演示图注被放在同一页，便于读者直接核对分面说明。真实项目应使用经过授权或用户提供的原图素材，并保留图注署名与许可信息。尺寸不足时，可适度缩放或调整栏宽，不以删除图注解决空间问题。\par
 模板仅记录一次已核验文章的可观察设计特征。即使文章出自同一期刊，不同年代、研究论文与综述、校正清样与正式版本也可能采用不同的页眉和引文形式。应以用户提供的实际PDF为准，不能仅凭期刊名称推断配色、字号或引用方式\SampleCiteOne。\par
 \subsection*{演示检查项}
 本样板用于展示可编译性、中文字形、标题层级、跨栏图注、页眉页脚和内部链接。图中没有实际数据，因此不应被当作科学插图引用。替换内容后，还需检查所有页的渲染结果，确认无截字、碰撞、异常空白和遗漏段落。\par
 \SampleBibliography
 \SampleSourceInfo
}


\usepackage{multicol,longtable,array,booktabs,tabularx}
\newcommand{\JournalName}{Nature Reviews Neuroscience}
% Source stroke RGB retained, rather than rounding to #41AD49.
\definecolor{NRGreen}{rgb}{.2539864182472229,.6789959669113159,.2859998345375061}
\definecolor{NRBlue}{HTML}{0068A5}
\definecolor{NRBox}{HTML}{F0F0F5}
\hypersetup{linkcolor=NRBlue,urlcolor=NRBlue,pdftitle={Nature Reviews Neuroscience风格中文中性样板},pdfauthor={原创中性示例；非期刊官方模板}}
\setlength{\columnsep}{12bp}\setlength{\parindent}{1em}\setlength{\parskip}{1.5bp}\setlength{\multicolsep}{6bp}
\renewcommand{\normalsize}{\fontsize{9.05bp}{12.1bp}\selectfont}
\ctexset{section={format=\sffamily\bfseries\fontsize{11bp}{13.5bp}\selectfont,beforeskip=8bp,afterskip=3bp},subsection={format=\sffamily\bfseries\fontsize{9.7bp}{12.1bp}\selectfont,beforeskip=6bp,afterskip=2bp}}
\newcommand{\NRCite}[1]{\textsuperscript{\hyperlink{ref:#1}{#1}}}
\newenvironment{NRColumns}{\begin{multicols}{2}\normalsize}{\end{multicols}}
\fancyhf{}
\AddToHook{shipout/foreground}{\begin{tikzpicture}[remember picture,overlay]
\draw[NRGreen,line width=3.5bp]([xshift=39.685bp,yshift=-35.765625bp]current page.north west)--([xshift=561.26bp,yshift=-35.765625bp]current page.north west);
\ifnum\value{page}=1
\node[anchor=base west,inner sep=0,font=\sffamily\fontsize{13.5bp}{17bp}\selectfont]at([xshift=39.685bp,yshift=-52.5bp]current page.north west){nature reviews neuroscience};
\node[anchor=base east,inner sep=0,text=NRBlue,font=\sffamily\fontsize{7.4bp}{9bp}\selectfont]at([xshift=561.26bp,yshift=-48.8bp]current page.north west){DOI：按目标来源核实};
\node[anchor=base west,inner sep=0,font=\sffamily\bfseries\fontsize{10bp}{12bp}\selectfont]at([xshift=39.685bp,yshift=-102.7108bp]current page.north west){Review article};
\draw[black,line width=.35bp]([xshift=39.685bp,yshift=-109.0905bp]current page.north west)--([xshift=561.26bp,yshift=-109.0905bp]current page.north west);
\else
\node[anchor=base west,inner sep=0,font=\sffamily\bfseries\fontsize{14bp}{17bp}\selectfont]at([xshift=39.685bp,yshift=-53.8bp]current page.north west){Review article};
\draw[black,line width=.35bp]([xshift=39.685bp,yshift=-109.1748bp]current page.north west)--([xshift=561.26bp,yshift=-109.1748bp]current page.north west);
\fi
\draw[black,line width=.35bp]([xshift=39.685bp,yshift=-754.4bp]current page.north west)--([xshift=561.26bp,yshift=-754.4bp]current page.north west);
\node[anchor=base west,inner sep=0,text=NRBlue,font=\sffamily\fontsize{7.4bp}{9bp}\selectfont]at([xshift=39.685bp,yshift=-769bp]current page.north west){Nature Reviews Neuroscience};
\node[anchor=base east,inner sep=0,font=\sffamily\fontsize{7bp}{9bp}\selectfont]at([xshift=561.26bp,yshift=-769bp]current page.north west){非官方中性样板\quad 中文译页\ \thepage};
\end{tikzpicture}}
\newcommand{\NeutralParagraph}{本段用于检查中文综述的连续阅读，不报告实际科学发现。正式译稿需要保留每个论点的对象、证据、比较方向和限定条件，分别核对正文、图注、表格和框文的完整范围。字体与行距可以适配中文，但不能以固定示例页数代替内容核对。\par}
\begin{document}\normalsize
% Regular-weight serif title intentionally follows the source's light hierarchy.
{\noindent\mdseries\fontsize{25bp}{29bp}\selectfont 中文综述题名占位：从完整材料\\到可追溯的研究论述\par}\vspace{35bp}
{\noindent\sffamily\bfseries\fontsize{9bp}{11.5bp}\selectfont 示例作者甲\textsuperscript{1,2}，示例作者乙\textsuperscript{1,3}，示例作者丙\textsuperscript{2,3}与示例作者丁\textsuperscript{1,*}\par}
\vspace{4bp}\hrule height .35bp\vspace{4bp}
\noindent\begin{minipage}[t][519bp][t]{388.0bp}\vspace{0pt}
{\sffamily\bfseries\fontsize{10bp}{12bp}\selectfont 摘要\par}\vspace{17bp}\hrule height .3bp\vspace{5bp}
{\fontsize{11.5bp}{16bp}\selectfont\setlength{\parindent}{0pt}
本页展示一套综述首页的结构，全部文字均为原创中性占位，不包含实际研究结果或临床建议。题名采用常规字重衬线体，作者区位于其下方，摘要与章节目录并列，以细线清楚区分。正式译稿必须完整替换题名、作者、单位、摘要和出版信息；本例的字数、段数与署名数量不构成期刊要求。\par\smallskip
综述中的概念、数据、证据与推论各有独立作用。翻译应沿用原作者的论证顺序，区分观察、关联、实验验证和未来设想，不把探索性框架写成已经建立的结论。原文中的否定、比较方向、不确定性及物种范围都需要保留；词句相近的段落也不能擅自合并。\par\smallskip
完整内容不只包括连续正文。图表、完整图注、框文、术语表、书目评注以及作者和出版声明，都有各自的来源位置和信息范围。重排中文时可以改变分页，原来的图号、文献号、术语对应和内容边界仍应可追溯。正文中仅提及、但未随文件提供的附件，不能算作已经取得或翻译。\par\smallskip
排版应保留来源的纸型、栏结构、层级和色彩关系。首页与普通正文页的题名、页眉及留白并不相同，不能用同一段固定高度替代实际测量。图像应在相应论述附近出现，长图注和中文表格按真实长度处理。所有文字仅用于检验本中性示例的阅读与版面，不对应任何疾病、实验或生物学主张。\par\smallskip
同一期刊的不同年代或文体，可能共享部分外观，同时采用不同的内容组织方式。应先检查所给首页、正文、图表页和末页，再决定哪些结构可以复用。完整性核对与视觉核对各有作用：前者确认材料没有漏译、重译或串接，后者检查图像边界、字形、层级与分页。两类核验通过以后，才形成可编辑源文件与最终阅读版本。\par}
\vfill\hrule height .3bp\vspace{4bp}
{\sffamily\fontsize{7.2bp}{9.4bp}\selectfont\noindent
\textsuperscript{1}作者单位一（原创中性占位）：机构、城市与国家按所给原文填写。
\textsuperscript{2}作者单位二（中性占位）。\textsuperscript{3}作者单位三（中性占位）。
*通信作者与联系方式按来源填写；本页无真实姓名、机构或邮箱。
仅供学术交流、出版物版权归原作者及出版商所有。\par}
\end{minipage}\hfill\vrule width .3bp\hfill
\begin{minipage}[t][519bp][t]{122.0bp}\vspace{0pt}\sffamily\fontsize{8.6bp}{11.5bp}\selectfont
{\bfseries\fontsize{10bp}{12bp}\selectfont 章节目录\par}\vspace{17bp}\hrule height .3bp\vspace{5bp}
引言\par\vspace{5bp}\hrule\vspace{5bp}
证据与概念的对应\par\vspace{5bp}\hrule\vspace{5bp}
图文与完整图注\par\vspace{5bp}\hrule\vspace{5bp}
研究汇总表与框文\par\vspace{5bp}\hrule\vspace{5bp}
术语、声明与书目\par\vspace{5bp}\hrule\vspace{5bp}
原文信息\par\vspace{5bp}\hrule
\end{minipage}
\clearpage
\begin{NRColumns}
\section*{正文与图文的衔接}
\NeutralParagraph
\subsection*{证据与描述范围}
不同材料可能反映不同层级的证据。原文将技术能力、观察发现和拟议机制分别陈述时，译文也应保持这一层次。一个示意图中的箭头不应被自动解释成已经获得因果验证\NRCite{1}。\par
\subsection*{完整段落的对应}
图像可能位于跨页续段之间，图注也可能出现在图像前一页。应先恢复原文的连续阅读顺序，再确定译文图文位置，避免丢失续句、重复图注或遗漏缩写。
\end{NRColumns}
\vspace{5bp}\noindent\begin{minipage}{\linewidth}\centering
\begin{tikzpicture}[font=\sffamily\fontsize{9bp}{12bp}\selectfont,node distance=20bp,
box/.style={draw=black!55,fill=black!3,minimum width=125bp,minimum height=60bp,align=center}]
\node[box](a){\textbf{A}\quad 来源材料\\原编号与内容范围};
\node[box,right=of a](b){\textbf{B}\quad 转写与翻译\\保留限定与比较};
\node[box,right=of b](c){\textbf{C}\quad 核验与交付\\完整正文及出处};
\draw[-{Stealth[length=2.2mm]},line width=.7bp](a)--(b);\draw[-{Stealth[length=2.2mm]},line width=.7bp](b)--(c);
\end{tikzpicture}\par\vspace{6bp}
{\sffamily\fontsize{7.6bp}{9.2bp}\selectfont\raggedright\textbf{图1｜材料转换的对应关系（原创中性占位，无实验数据）。}A，输入材料按实际所给版本记录。B，转写与翻译分别保留原编号、概念和限定。C，输出需要核对正文、图注、表格、框文和文末信息的完整性。箭头只示意阅读顺序，不表示因果关系、效应方向、时间长度或数量变化。本图没有样本、数据点、误差线、统计检验或比例尺。\par}
\end{minipage}\vspace{7bp}
\begin{NRColumns}
\subsection*{科学原图与标签释义}
保留原图时应完整翻译图注，并提供非专有英文标签、整句说明、图例和流程节点的中文对应。基因、蛋白、数字及科学坐标不能因排版需要而改写。长标签释义可独立附于文末，并清楚标明图号\NRCite{2}。\par
\NeutralParagraph
\subsection*{裁图与颜色核验}
图像边界应逐幅测量，不沿用其他文章的固定坐标。处理下载来源时须定位实际对象，确认科学图的全部面板、边缘和尺度均被保留。只改变裁区并不能证明裁区外的文字已从文件中移除。\par
对照来源渲染图检查文字和像素，可以识别通用清理造成的字体宽度、面板位置或颜色变化。检查记录应区分科学内容变化与纯版式适配。
\end{NRColumns}
\clearpage
\section*{表1｜汇总材料的独立记录（原创占位）}
{\sffamily\fontsize{8bp}{10.5bp}\selectfont\setlength{\tabcolsep}{5bp}
\begin{longtable}{@{}p{70bp}p{105bp}p{\dimexpr\linewidth-195bp\relax}@{}}
\toprule 材料类别 & 需要核对的字段 & 中性记录说明\\\midrule\endfirsthead
\multicolumn{3}{l}{\bfseries 表1（续）｜汇总材料的独立记录}\\[4bp]
\toprule 材料类别 & 需要核对的字段 & 中性记录说明\\\midrule\endhead
正文 & 论点、对象与限定 & 对应来源的完整段落；表中摘要不代替正文。\\[6bp]
图像 & 面板、图例与尺度 & 图形全部为演示占位；实际图像按目标来源核对。\\[6bp]
图注 & 样本、比较与缩写 & 实际源文报告的数值和统计限定不因图中已有标记而省略。\\[6bp]
框文 & 所有标题与条目 & 框文独立于正文维护，不把中性示例改写成研究发现。\\[6bp]
术语 & 缩写、定义与范围 & 定义需依照当前文章，不能从另一篇文章直接继承。\\[6bp]
书目 & 原文条目与编号 & 保留英文书目，中文评注另行呈现；链接对应原编号。\\[6bp]
声明 & 贡献、利益与许可 & 来源实际提供的每项声明均应完整覆盖；缺失字段不编造。\\[6bp]
\bottomrule\end{longtable}
\noindent 表注：本表无真实论文、队列、平台性能或实验数据。所有条目只记录材料层级，不是正式研究的固定纳入标准。\par}
\vspace{8bp}\noindent\colorbox{NRBox}{\begin{minipage}{\dimexpr\linewidth-2\fboxsep\relax}
\sffamily\fontsize{8.5bp}{11.5bp}\selectfont
\textbf{框1｜完整内容的覆盖边界（原创中性占位）}\par\smallskip
\textbf{实际取得的材料。}正文、图表、框文、术语表、作者声明和原文信息需分别登记。原文仅引用某个独立补充文件，并不意味着该文件已随输入提供。\par\smallskip
\textbf{翻译与摘要的区别。}全文转换应保留全部论述、定义、条件和编号。中性样板只用于演示结构，不能用模板的短文字代替实际文章的完整内容。\par\smallskip
\textbf{版式与科学内容。}字体回退、中文行距和分页可以调整；作者给出的比较方向、因果强度、物种及适用范围应保持不变。\par\smallskip
\textbf{发现来源差异时。}可见页面、底层文本和出版字段出现差异时，先回到原页核对，再明确说明依据；不从文件名推填缺失的在线日期、卷期或页码。\par
\end{minipage}}\par\vspace{8bp}
\begin{NRColumns}\subsection*{表格与框文的复核}
\NeutralParagraph
多页表格可使用重复列头，保持分组与表注归属。框文可在清楚的边界处分段续排，但不能截去末尾项目。字体和行距的紧凑处理不能替代对每一行的完整翻译。
\end{NRColumns}
\clearpage
\noindent\begin{minipage}[t]{254.7bp}\vspace{0pt}
\section*{术语表（原创中性占位）}
{\sffamily\fontsize{8.1bp}{10.7bp}\selectfont
\textbf{来源版本}\par 实际提供的文件及其出版身份，不用文件名推断后续版本。\par\smallskip
\textbf{内容覆盖}\par 能逐项映射到输出的实际材料范围，不包括未取得的独立附件。\par\smallskip
\textbf{图注}\par 解释原图面板、限定、缩写及统计信息的独立文字，不是可省略的装饰。\par\smallskip
\textbf{术语对应}\par 当前文章中概念、缩写与定义的可追溯关系。\par\smallskip
\textbf{书目目标}\par 供原编号引文跳转的对应条目。\par\smallskip
\textbf{版式适配}\par 因中文长度和字体回退作出的阅读调整，不改变科学内容。\par\smallskip
\textbf{原文信息}\par 文末对作者、英文原题、期刊、日期、卷期页码和DOI的完整来源记录。\par}
\section*{核验记录}
\NeutralParagraph
逐段检查与逐页渲染应分别保留记录。文本提取完整不等于图像边界完整，页面看起来整齐也不证明内部引文链接都有效。\par
\end{minipage}\hfill\begin{minipage}[t]{254.7bp}\vspace{0pt}
\section*{声明（原创占位）}
本样板无真实研究、实验对象、资助、作者贡献或利益冲突。实际文章中的这些声明应按源文分别处理，不因它们位于书目后而遗漏。\par
本例以普通文字刊名与可移植字体展示已核验结构，不含期刊官方标识图像、真实论文题名、作者、DOI或原图，不代表期刊官方模板或出版认可。\par
\SampleBibliography
\SampleSourceInfo
仅供学术交流、出版物版权归原作者及出版商所有。\par
\end{minipage}
\end{document}
